% results.tex -- Results

We ran SA on all eight Pfam families (Table~\ref{tab:families}) and observed that the method consistently generated novel protein sequences while preserving family-level amino acid statistics, across a ten-fold range of family sizes and a seven-fold range of sequence lengths (Fig.~\ref{fig:cross-family}; SI Appendix, Table~\ref{tab:results}). We first assessed sampler accuracy by comparing the outputs of the unadjusted Langevin algorithm (ULA) and its Metropolis-adjusted variant (MALA): acceptance rates exceeded 99.6\% in every family. The Metropolis correction rejected fewer than 0.4\% of proposed moves at this step size (SI Appendix, Table~\ref{tab:sampling-diag}). Thus, the unadjusted and adjusted updates rarely differed in this regime, and we adopted the computationally cheaper ULA in all subsequent experiments. Generated sequences were not near-duplicates of stored patterns: zero sequences across all eight families exceeded 90\% identity to any stored sequence (SI Appendix, Table~\ref{tab:identity-dist}), and the generated ensemble was insensitive to initialization strategy (SI Appendix, Table~\ref{tab:random-init}). In the generation regime, the sampler operated at inverse temperature $\beta \approx 2\beta^*$, where the critical inverse temperature $\beta^*$ is the value at which the energy landscape changes character: below $\beta^*$ the stored patterns blur into a single broad well, while above it the landscape resolves into distinct basins around individual patterns (SI Appendix, Section~\ref{si:beta-star-theory}). At $\beta \approx 2\beta^*$, the generated amino acid composition closely matched each natural family's: the Kullback--Leibler (KL) divergence between the generated and natural amino acid compositions, which falls to zero only when two distributions are identical, stayed at or below $0.06$ in every family and below $0.02$ in three of them. At the same time, novelty ranged from $0.40$ (Defensin\_beta) to $0.65$ (WW), indicating departure from any single stored memory. All generated sequences were distinct in every family, with mean pairwise sequence identity of $35$--$65\%$ (equivalently, mean pairwise Hamming distance $0.36$--$0.65$; SI Appendix, Table~\ref{tab:nn-band}). Thus, the sampler produced a diverse ensemble rather than collapsing to a single mode. Nearest sequence identity to a stored pattern was $51$--$66\%$, within each family's empirical within-family nearest-neighbor identity envelope (each natural sequence's identity to its closest neighbor in the same family; SI Appendix, Table~\ref{tab:nn-band}) in seven of eight families. In Pkinase, generated identity ($0.515$) exceeded the maximum nearest-neighbor identity among natural Pkinase sequences ($0.454$), consistent with that family's consensus-proximal generation. Together, generated sequences remained within the observed sequence-similarity envelope while departing from stored patterns. By contrast, bootstrap samples had zero novelty by construction and Gaussian perturbation produced negligible novelty ($<0.02$), indicating that small additive noise is insufficient to escape the basin of the nearest stored pattern.

Position-specific analysis of the Kunitz domain revealed why SA preserves family statistics (SI Appendix, Fig.~\ref{fig:sequence-analysis}). Sequences generated by stochastic attention preserved all 11 strongly conserved positions of the 53-column seed alignment (those exceeding 90\% conservation), including all six cysteines forming the three disulfide bonds essential to the Kunitz fold, while introducing 17--25 substitutions concentrated at the variable positions. Per-position Shannon entropy correlated between the SA-generated ensemble and the stored alignment (Pearson $r{=}0.968$). Thus, the method identified which positions tolerate variation and which do not. This behavior followed from the Boltzmann sampling mechanism: at strongly conserved positions, all stored patterns agreed and the softmax weights concentrated probability on the consensus residue; at variable positions, stored patterns diverged and the sampler drew from the local distribution of tolerated amino acids.

Subsampling the WW domain at $K \in \{20, 50, 100, 200, 400\}$ with five independent replicates per condition characterized how generation quality changed with decreasing family size (SI Appendix, Fig.~\ref{fig:scaling}). Stochastic attention maintained low KL divergence across the entire range ($0.019 \pm 0.008$ at $K{=}20$, improving to $0.010 \pm 0.002$ at $K{=}400$), while the PCA dimensionality grew from $d{=}17$--$18$ components at $K{=}20$ to $d{=}182$--$183$ at $K{=}400$ and the critical inverse temperature tracked this growth ($\beta^*{=}2.7$ to $5.5$). Despite the roughly ten-fold increase in effective dimensionality, compositional fidelity remained stable, suggesting that the entropy inflection criterion adapted the operating temperature to the geometry of the memory matrix at each scale. Novelty increased with $K$ (from $0.47$ to $0.74$), reflecting the richer landscape of a larger memory set. Nearest sequence identity decreased correspondingly (from $0.71$ to $0.56$), indicating that the sampler explored further from stored patterns when more patterns were available to collectively shape the energy landscape. Replication across three additional families at $K{=}20$ (SH3, Kunitz, zf-C2H2) confirmed that this robustness is not family-specific (SI Appendix, Table~\ref{tab:multifamily-k20}).

The scaling study revealed a consistent pattern in how the sampler self-tunes: the more independent variation a family contained, the higher the inverse temperature the sampler selected for it. The critical inverse temperature $\beta^*$ tracked the family's PCA dimensionality $d$, the number of principal components needed to capture 95\% of its sequence variation. This prompted us to ask whether $\beta^*$ could be predicted from alignment statistics alone, without a temperature sweep (Fig.~\ref{fig:beta-prediction}). A two-parameter model linear in $\sqrt{d}$ predicted $\beta^*$ from PCA dimensionality alone. Fitted on the eight independent Pfam families, it accounted for $95\%$ of the variance in $\beta^*$ ($R^2{=}0.95$) and generalized under leave-one-family-out cross-validation (SI Appendix, Section~\ref{si:beta-star-theory}). Pooling all 33 data points (the 8 families and 25 WW scaling replicates) gave a nearly identical fit, $\beta^* \approx 1.52 + 0.28\sqrt{d}$ (bootstrap standard error [SE]: intercept $\pm 0.07$, slope $\pm 0.007$; $R^2{=}0.97$). The square-root dependence on $d$ follows from concentration of measure on the unit sphere $\mathbb{S}^{d-1}$: the similarity between a random probe and the stored patterns scales as $\mathcal{O}(1/\sqrt{d})$, so the inverse temperature needed to resolve those similarities grows as $\sqrt{d}$. The fitted slope, however, was smaller than geometry alone predicts: $0.28$, against the ${\approx}1.6$ expected for random patterns under a Gaussian random-score model. That gap is a signature of real sequence structure. Each stored pattern has unit inner product with itself but lower similarity to the other patterns (a gap of $0.55$--$0.82$ to its next-nearest pattern), so when the patterns themselves serve as probes the softmax concentrates on the self-pattern at a lower $\beta$ than the random-score model predicts (SI Appendix, Section~\ref{si:beta-star-theory}). Adding further predictors (mean column entropy, effective number of sequences, spectral concentration) produced negligible improvement ($\Delta R^2 < 0.001$), indicating that PCA dimensionality captured most of the relevant variation. This regression provides a practical shortcut: for a new family, $\beta^*$ can be estimated from its PCA dimensionality and used to set the generation temperature without a sweep.

Having characterized how SA self-tunes in isolation, we next compared it against three established learned baselines, profile HMM emit, EvoDiff (MSA-conditioned diffusion), and the MSA Transformer (Gibbs masked language model sampling), across all eight families (Fig.~\ref{fig:baseline-comparison}); a physics-based Potts model was compared separately. Profile HMM emit and the MSA Transformer achieved high novelty scores, but their nearest sequence identity to stored patterns fell to $9$--$59\%$ across the eight families. Because the identity reported for each generated sequence is a nearest-neighbor quantity, the comparable natural quantity is the within-family nearest-neighbor identity, each natural sequence's identity to its closest neighbor in the same family. By this measure, the drifted baselines fell below the lowest nearest-neighbor identity of any natural sequence: profile HMM emit reached $11$--$21\%$ in Pkinase, PDZ, and RRM, whereas no natural sequence in those families is less than $22$--$28\%$ identical to its nearest neighbor (SI Appendix, Table~\ref{tab:nn-band}). Each family's natural within-family nearest-neighbor identity envelope is shown in green (Fig.~\ref{fig:baseline-comparison}C); a sequence falling below its family's box may have moved outside the family, whereas one rising above it is a near-copy of stored patterns with insufficient novelty. EvoDiff balanced the metrics more effectively in shorter families but required pretraining on millions of UniRef50 MSAs and scaled poorly with sequence length (${\sim}14$ hours for 150 Pkinase sequences on CPU, though GPU inference would reduce this, versus seconds for SA on the same hardware; SI Appendix, Table~\ref{tab:timing}). Taken together, none of the three learned baselines combined compositional fidelity, novelty, and family-consistent sequence identity: the profile HMM and MSA Transformer fell outside the family's identity range, while EvoDiff achieved the balance only in shorter families and only with large-scale pretraining.

Stochastic attention simultaneously achieved compositional fidelity, novelty, and family-consistent sequence identity (Fig.~\ref{fig:baseline-comparison}). It held KL divergence below $0.05$ in seven of eight families. Its novelty of $0.40$--$0.65$ indicated departure from stored patterns. Sequence identity stayed at $51$--$66\%$, within each family's empirical within-family nearest-neighbor identity envelope for seven of eight families, with Pkinase marginally above (SI Appendix, Table~\ref{tab:nn-band}). This combination (compositionally faithful, novel, and within the family's sequence-similarity envelope) distinguished SA from the baselines tested (which differ from SA in architecture and training requirements). A physics-based model that fits that structure explicitly gives a sharper test. We therefore compared SA against a Potts model fitted via pseudo-likelihood maximization (plmDCA) across all eight families (SI Appendix, Table~\ref{tab:potts}). Potts models explicitly fit pairwise coupling parameters $J_{ij}(a,b)$ between all position pairs, giving them direct access to coevolutionary signals. The comparison between the Potts model and SA revealed comparable compositional fidelity and output quality in six of eight families, with SA achieving higher novelty in five of eight. However, the Potts model showed limitations in the most data-starved families. In Pkinase its novelty was $0.22$, against $0.48$ for SA, showing the Gibbs sampler remained near stored patterns. In Defensin\_beta the Potts method KL of $0.112$ was the highest of any method (consistent with its high parameter count). The number of coupling parameters scales as $L(L{-}1)q^2/2$ for alignment length $L$ and alphabet size $q{=}20$, yielding data-to-parameter ratios as low as $2.5 \times 10^{-6}$. That SA achieved comparable or better output without fitting any parameters suggests that the Hopfield energy implicitly captures some of the correlational structure that pairwise coevolutionary models fit explicitly, but through the softmax attention operation rather than through parameter estimation.

Pairwise mutual information (MI) analysis tested whether SA preserved residue covariation, which reflects structural contacts and functional couplings that single-site statistics cannot capture~\cite{morcosDirectCouplingAnalysis2011} (SI Appendix, Table~\ref{tab:mi}; Figs.~\ref{fig:mi-scatter}--\ref{fig:mi-barplot}). Pearson correlations between stored and SA-generated MI vectors ranged from $r{=}0.53$ to $0.92$ across seven of eight families, with Pkinase ($r{=}0.05$) the lone exception. The Pkinase alignment was shallow, only $K{=}37$ sequences over $L{=}262$ positions, which could in principle make its covariation hard to estimate. However, this was likely not the case because the Potts model recovered the covariation on the identical alignment ($r{=}0.95$). The failure lay instead in SA generation, which in Pkinase stayed closer to the family consensus than in any other family (mean identity to the consensus $0.73$, the highest observed; SI Appendix, Section~\ref{si:consensus-reg}). Because near-consensus sequences vary little from one another, they carried almost none of the position-to-position covariation that MI measures. Thus, stochastic attention preserved covariation in the compact families but not in the longest, most data-starved one. Further, stochastic attention exceeded profile HMMs (whose MI correlations spanned $r{=}0.07$--$0.83$) in seven of eight families. Ranked by covariation fidelity, the four methods followed the order HMM $<$ SA $<$ Potts $\leq$ EvoDiff in three of eight families. This order tracked how explicitly each method represents covariation: profile HMMs model only position-independent frequencies, stochastic attention captures couplings implicitly through the Hopfield energy, Potts models fit pairwise couplings explicitly, and EvoDiff adds deep coevolutionary priors from large-scale pretraining. That stochastic attention outranked the position-independent baseline suggested the Hopfield energy captured correlational structure single-site models cannot encode, without fitting any pairwise parameters.

We assessed structural plausibility of SA-generated sequences with two independent structure predictors, ESMFold~\cite{linEvolutionaryScaleModeling2023} and AlphaFold2~\cite{jumperHighlyAccurateProtein2021}. We predicted folded structures for 50 sequences per method per family (Fig.~\ref{fig:structure-validation}). We compared SA-generated and stored natural sequences on two structural metrics. For each metric, a Wilcoxon rank-sum test assessed significance and Cohen's $d$ measured the effect size. The first structural metric, the predicted local distance difference test (pLDDT) score, is the predictor's own per-residue confidence that each position was placed correctly, reported from 0 to 100. Under ESMFold, sequences generated by stochastic attention achieved mean pLDDT scores statistically indistinguishable from stored sequences in six of eight families (no significant difference, $p > 0.05$). The pLDDT was significantly higher in Kunitz ($90.8$ vs.\ $89.3$, $p{=}0.001$, $d{=}0.67$). At least 86\% of SA-generated sequences exceeded the pLDDT${=}70$ confidence threshold in seven of eight families, with SH3 the lone exception ($54\%$). The second structural metric, the template modeling (TM) score, ranges from 0 to 1 and rises as a predicted structure superimposes more closely on the reference. In six of eight families, the predicted structures of SA-generated sequences matched the reference significantly better than stored sequences ($p < 0.05$; $d > 1.0$ in RRM, WW, PDZ, and Pkinase). To test whether these elevated TM-scores reflected a predictor preference for consensus-like sequences, we asked whether predicted structural quality tracks consensus proximity within each family. The correlation was weak and inconsistent (per-family Spearman $\rho$ between TM-score and identity to the consensus ranged from $-0.13$ to $0.53$, pooled $\rho{=}0.17$; pLDDT pooled $\rho{=}0.23$; SI Appendix, Section~\ref{si:consensus-reg}). Thus, the elevated TM-scores were likely not an artifact of consensus proximity. As an independent cross-validation, the AlphaFold2 predictions corroborated the ESMFold findings while resolving the single negative result (SI Appendix, Figs.~\ref{fig:structure-esmfold-af2}--\ref{fig:esmfold-af2-concordance}). For SH3, the only family where ESMFold had indicated a pLDDT deficit, AlphaFold2 showed no pLDDT difference ($p{=}0.99$) and significantly higher TM-scores for SA generation ($0.736$ vs.\ $0.721$, $p < 0.001$). For the highest-confidence representative from each family, the AlphaFold2 prediction aligned to the experimentally determined reference with sub-angstrom root-mean-square deviation (Fig.~\ref{fig:structure-gallery}). The concordance between two independent structure predictors across all eight families strengthens the computational evidence that SA-generated sequences are predicted to adopt the target-family fold.

As a complementary, sequence-level assessment of biological plausibility, we scored all SA-generated and stored sequences using ESM2-650M~\cite{linEvolutionaryScaleModeling2023}, a protein language model pretrained on UniRef50 sequences that is entirely independent of SA (SI Appendix, Table~\ref{tab:esm2}). Overall, the model scored SA-generated sequences comparably to natural family members, with mean pseudo-perplexity $5.71$ versus $5.58$ for stored sequences. Lower pseudo-perplexity indicates a sequence the model finds more typical of natural proteins. In five of eight families, SA scored as well as the stored sequences or better: indistinguishable in four (RRM, zf-C2H2, Pkinase, Defensin\_beta; Wilcoxon $p > 0.29$) and significantly \emph{better} in Kunitz ($4.47$ vs.\ $4.98$, $p < 0.001$, $d{=}{-}0.74$). In the other three families (SH3, WW, PDZ), SA scored modestly higher ($d{=}0.60$--$1.03$), though all generated sequences remained within the range of natural proteins. Beyond this language-model scoring, we cross-referenced SA-generated substitutions against published deep mutational scanning (DMS) datasets to ask whether the Boltzmann sampling procedure preferentially generates individually tolerated mutations (SI Appendix, Tables~\ref{tab:esm2} and \ref{tab:dms}). For the PDZ domain (DLG4/PSD-95 PDZ3~\cite{mclaughlinSpatialArchitecture2012}), $91.7\%$ of $4{,}313$ matched substitutions were classified as experimentally tolerated, and for the SH3 domain (GRB2 SH3~\cite{faureDoubleDMS2022}), $90.4\%$ of $2{,}396$ matched substitutions were tolerated. Because SA preferentially substitutes at variable, more tolerant positions, comparing these rates against the overall dataset tolerance rate ($77.5\%$ and $67.5\%$) conflates position selection with residue choice. We therefore used a position-matched null that asks, for each position SA substituted, how often a randomly chosen tested substitution at that same position is tolerated. Against this stricter null, SA-generated substitutions remained significantly enriched in both families, indicating that SA selects function-compatible residues beyond simply targeting tolerant positions. The enrichment held for PDZ ($91.7\%$ vs.\ $84.3\%$, $1.09\times$, $p < 10^{-77}$) and SH3 ($90.4\%$ vs.\ $74.9\%$, $1.21\times$, $p < 10^{-121}$), and a conservation-matched null gave the same result (SI Appendix, Table~\ref{tab:dms-null}). The continuous DMS fitness scores of SA-generated substitutions were likewise significantly higher than those of random single mutations (Wilcoxon $p < 10^{-20}$; SI Appendix, Tables~\ref{tab:esm2} and \ref{tab:dms}). Taken together, these results indicate that the Boltzmann sampling procedure preferentially generates substitutions that are individually compatible with function in DMS assays, choosing both tolerant positions and tolerated residues, without any explicit fitness objective. However, they do not by themselves establish that every full generated sequence is functional, and the DMS evidence is limited to these two domains.

Finally, the beta defensin family (PF00711) provided an opportunity to test whether SA generation preserved not only sequence-level statistics but also the biophysical properties that underpin antimicrobial function~\cite{lazzaroAntimicrobialPeptides2020} (SI Appendix, Fig.~\ref{fig:defensin-biophysics}). Sequences generated by stochastic attention preserved the six-cysteine disulfide scaffold ($6.01 \pm 0.08$ mean cysteines, compared to $6.02 \pm 0.15$ for natural sequences). They maintained a cationic charge profile ($+6.1 \pm 1.6$), with the narrowest charge distribution of any method. They also occupied the same region of charge-hydrophobicity space as natural defensins. By a minimal antimicrobial plausibility filter (net charge $\geq +2$, hydrophobic fraction in $[0.3, 0.7]$, $\geq 4$ cysteines), $51\%$ of SA-generated sequences passed, compared to $42\%$ of stored sequences, $25\%$ of HMM-emitted sequences, and $16\%$ of MSA Transformer sequences. This implicit enforcement of biophysical constraints arose because the energy landscape concentrated probability near the stored patterns, which already satisfied these constraints by virtue of being functional proteins. No explicit objective for charge, hydrophobicity, or disulfide topology was required.
